% ============================================================
% Chapter 4: Proposed Method: Specification and Design
%   Section: The Fact-Checking Pipeline
%     Subsection: Evidence Ranking and the Pertinence Gate
% Included from chapters/04-proposed-method/06-fact-checking-pipeline/fact-checking-pipeline.tex
% Source: docs/decisions/retrieval.md; ranking/ranking_retrieval.py;
% settings.py (RANKING_*, EVIDENCE_MIN_PERTINENCE).
% ============================================================

\subsection{Evidence Ranking and the Pertinence Gate}
\label{subsec:ranking}

The fetched evidence is re-ranked by a score of four independent
factors:
\[
  \textit{relevance} = 0.45\cdot\textit{semantic}
    + 0.20\cdot\textit{lexical}
    + 0.20\cdot\textit{recency}
    + 0.15\cdot\textit{reliability}.
\]

\begin{itemize}
    \item \textbf{Semantic}: the embedding similarity between the claim
    and the evidence.
    \item \textbf{Lexical}: the share of the claim's own terms the
    evidence contains, anchors counted double. It is the counterweight
    to the semantic factor, which scores a page \emph{about the same
    subject} as highly as a page that addresses the claim --- exactly
    how the ACME etymology pages won.
    \item \textbf{Recency}: an exponential decay with a half-life of a
    year, so old evidence loses weight gradually rather than being
    discarded.
    \item \textbf{Reliability}: a rating per domain, from the source
    catalogue. Only the 37 configured domains have a real rating; any
    other domain gets a neutral 0.5 and is flagged as unrated, so the
    default is never presented as a rating.
\end{itemize}

The weights must sum to one, and the backend refuses to start if they
do not: the configuration template once shipped weights summing to 1.2,
and nothing noticed. The best five sources per claim are kept, at most
one per domain.

\paragraph{The pertinence gate.} In the same step, each source gets a
second score that leaves recency and reliability out:
\[
  \textit{pertinence} = 0.5\cdot\textit{semantic} + 0.5\cdot\textit{lexical}.
\]
A recent, reliable page about something else is still about something
else, and letting those two factors lift it would repeat the mistake
the gate exists to stop. A source whose pertinence is below a threshold
(0.25 by default) is cut \emph{before the model sees it}, with a reason
that names the figure: it cannot be cited, counted towards the evidence
minimum or treated as corroboration. The threshold is deliberately low
--- a floor against off-target retrieval, not a relevance ranking ---
because cutting real evidence turns a checkable claim into
\texttt{UNVERIFIED}, another way of being useless. It is a per-run
threshold, unlike the weights, because it is the knob that trades a
confident wrong answer for an honest one. Its default was reasoned, not
fitted: no labelled set existed when it was set.
